\clearpage
\section{Results Comparison on Metallicity}

The metallicity task (Materials Project, \MetalN{} compositions, \MetalPos{}\% metals) was
trained with 1,000 labeled compositions and five semi-supervised methods, the masked-autoencoder
features and nine supervised baselines (\MetalInRuns{} models per unlabeled pool). Every score
is the mean over 3 seeds on the same \MetalInNTest{} held-out test compositions.

\subsection*{Comparison of SSL Methods}

\resfig{\textwidth}{results/metallicity/indist/metallicity_methods_comparison_balanced.png}{Comparison
of SSL methods on the Metallicity dataset (unlabeled pool: same database, labels hidden).}

\maybeinput{table_metallicity_indist}
\maybeinput{table_metallicity_gnome}

With the unlabeled compositions taken from the same database, \MetalInBestMethod{}
(\texttt{\MetalInBestConfig}) reaches a composite score of \MetalInBestComp{}, against
\MetalInBaseComp{} for the Optuna-tuned \MetalInBaseConfig{} baseline (\MetalInImpr\%).
Self-Training, Cluster-then-Label and the self-supervised features are statistically tied with
the baseline (differences below one standard deviation). The graph methods (Label Propagation
and Label Spreading) are clearly behind (composite $\approx 0.85$): the $k$-NN graph built on
142 descriptors does not separate metals from insulators as well as a boosted tree does.

With GNoME as the unlabeled pool, no SSL method beats the baseline: the best one,
\MetalGnBestSSL{}, reaches \MetalGnBestSSLComp{} (\MetalGnImpr\%), and the graph methods fall
to 0.75--0.79. GNoME compositions follow a different distribution (Figure~\ref{fig:shift}), so the labels
that are propagated or pseudo-labelled on them are less reliable.

\subsection*{Effect of the Number of Unlabeled Compositions}

\resfig{\textwidth}{results/metallicity/indist/metallicity_learning_curves.png}{Composite score
and F1-score versus the number of unlabeled compositions (metallicity, same-database pool).}

\maybeinput{poolsize_metallicity_indist}
\maybeinput{poolsize_metallicity_gnome}

Adding more unlabeled compositions does \emph{not} help: every method is best with the
smallest pool (1,000 unlabeled) and degrades steadily up to the largest one (60,945
for the same database, 100,000 for GNoME). Self-Training is the most
robust, because its 0.9 confidence threshold only accepts pseudo-labels the model is sure
about. The graph methods degrade fastest, because each additional unlabeled node also spreads
errors.

\subsection*{Label Budget and Supervised Model Zoo}

\resfig{0.8\textwidth}{results/metallicity/indist/metallicity_label_budget_curve.png}{Label-budget
curve for metallicity: best supervised baseline versus best SSL configuration.}

\maybeinput{budget_metallicity_indist}
\maybeinput{zoo_metallicity_indist}

Labels matter much more than unlabeled data: going from 1,000 to 5,000 labels raises the
composite from about 0.87 to 0.90, while using all \MetalN{} compositions brings the fully
supervised model to \MetalInFullComp{} (ROC-AUC \MetalInFullAUC{}). The SSL and baseline curves
overlap at every budget (differences of at most 0.001). Among supervised
models, gradient boosting (LightGBM, XGBoost) and Random Forest are equivalent; linear, $k$-NN
and MLP models are about 0.02 lower.

\subsection*{Selected Model}

The deployed model keeps the architecture selected at 1,000 labels and is re-trained on all
training labels plus the unlabeled pool. On the held-out test set it reaches a composite score
of \MetalSelComp{} (ROC-AUC \MetalSelAUC{}, F1 \MetalSelFone{}, accuracy \MetalSelAcc{}).

Consequently, \textbf{\MetalSelMethod{}} (\texttt{\MetalSelConfig}) is selected as the primary
architecture for the deployment phase of the metallicity prediction module.
